% Guia do professor — Capítulo 9: Boletins e Análise de Cartas
\section{Capítulo 9 --- Boletins Meteorológicos e Análise de Cartas}

\subsection*{Visão geral e ritmo}

Capítulo de alfabetização operacional: METAR, modelo de estação, redução ao
nível do mar e análise (manual e objetiva). \textbf{Ritmo: 1 aula
expositiva + 1 aula prática} com carta impressa e lápis — a prática é o
coração.

\begin{itemize}
  \item \textbf{Aula 1} — METAR grupo a grupo (com boletins do dia!),
        modelo de estação, redução ao nível do mar. Casos~1--2 (10\,min).
  \item \textbf{Aula 2 (prática)} — cada dupla analisa a mesma carta de
        superfície impressa: isóbaras de 4 em 4\,hPa, centros, frentes
        tentativas. Fechar comparando com a análise oficial e com o
        Caso~3.
\end{itemize}

\subsection*{Objetivos de aprendizagem}

\begin{enumerate}
  \item decodificar METARs sem tabela de consulta (grupos principais);
  \item plotar e ler o modelo de estação (barbela, pressão codificada);
  \item explicar e aplicar a redução ao nível do mar;
  \item traçar isóbaras seguindo as regras e identificar centros;
  \item descrever a análise objetiva (Cressman/Barnes) e seu papel
        histórico.
\end{enumerate}

\subsection*{Pré-requisitos a reativar}

Hipsométrica (Cap.~1); oktas (Cap.~6); $T_d$/UR (Cap.~4); convenção de
direção do vento (Cap.~1 — o erro clássico reaparece na barbela).

\subsection*{Armadilhas conceituais frequentes}

\begin{itemize}
  \item \textbf{Barbela ao contrário}: a haste aponta de onde o vento vem;
        metade da turma erra na primeira vez.
  \item \textbf{Pressão codificada}: 052 $\to$ 1005{,}2 (não 905{,}2);
        a regra do valor mais próximo de 1000 e sua exceção em furacões.
  \item \textbf{QNH × QFE × pressão da estação}: três pressões diferentes;
        cartas usam a reduzida.
  \item \textbf{Isóbaras ``ligando pontos''}: interpolação proporcional,
        não conexão de estações.
  \item \textbf{Análise objetiva como verdade}: é filtro; o raio de
        influência decide o que sobrevive (T5/N3).
\end{itemize}

\subsection*{Demonstração de baixo custo}

METARs ao vivo do REDEMET/aviationweather.gov no projetor: decodificar o
aeroporto da cidade em tempo real. Para a prática: imprimir a carta de
superfície da Marinha/CPTEC do dia anterior sem análise.

\subsection*{Problemas sugeridos do livro (exercícios do Cap.~9)}

Aquecimento: decodificação de boletins; plotagem de estações.
Aprofundamento: análise de mini-cartas com 10--15 estações. Síntese:
``que observações faltam sobre os oceanos e como os satélites (Cap.~8)
preenchem?''. \emph{Confira a numeração na sua cópia (v1.02b).}

\subsection*{Uso do notebook \texttt{cap09\_cartas.ipynb}}

\begin{center}
\begin{tabular}{lll}
\hline
Caso & Conteúdo & Momento sugerido \\
\hline
1 & Parser de METAR (Brasil) & Aula 1, após o código \\
2 & Modelo de estação (MetPy) & Aula 1, após o modelo \\
3 & Cressman: pontos $\to$ isóbaras & Aula 2, após a prática manual \\
4 & Redução ao nível do mar: antes × depois & Aula 1 ou 2 \\
\hline
\end{tabular}
\end{center}

Momento pedagógico forte: mostrar o Caso~3 \emph{depois} da análise manual
— os alunos reconhecem no algoritmo o que a mão deles acabou de fazer.

\subsection*{Exercícios complementares e soluções}

Nas notas: T1--T5 e N1--N4 (o par N1$\to$N2 constrói um mini-sistema de
plotagem). Soluções em \texttt{cap09\_solucoes.ipynb}
(\emph{não distribuir}); N3 traz validação cruzada de verdade — primeira
aparição de um conceito de aprendizado estatístico no curso.

\subsection*{Conexões com o resto do curso}

Hipsométrica $\to$ Cap.~1; oktas/teto $\to$ Cap.~6; isóbaras$\to$vento
$\to$ Cap.~10; frentes $\to$ Cap.~12; análise objetiva $\to$ assimilação
(Cap.~20).
